\documentclass[11pt]{article}
\usepackage{preamble}
\renewcommand{\answer}[1]{}

\begin{document}

\ladderheading{AP Statistics \textperiodcentered\ Topic 4.1 (CED 2.3) \textperiodcentered\ B: 2026-10-19 \textperiodcentered\ E: 2026-11-02}{Can Chance Make a Streak?}

\focusbanner{TODAY'S PROBLEM}

Suppose a coin advertised as fair is flipped 10 times and gives HHHHHHTTTT. To investigate what chance alone could produce, a class simulates 100 sets of 10 independent fair-coin flips. For each set, the class records the longest run of matching outcomes; the dotplot shows the results.\par\smallskip \textbf{Claim A:} ``The coin is rigged---six heads in a row cannot be chance.''\par \textbf{Claim B:} ``Streaks like this can happen by chance; simulation can tell us how unusual this one is.''\par
\begin{center}
\begin{tikzpicture}
\begin{axis}[
  width=0.85\linewidth,
  height=1.60in,
  ymin=0, ymax=33, ytick=\empty, axis y line=none, axis x line=bottom,
  xlabel={Longest run in one set of 10 fair-coin flips}, tick label style={font=\small}, label style={font=\small}
]
\addplot[only marks, mark=*, mark size=2.4pt, color=dayblue] coordinates {
  (2,1)
  (2,2)
  (2,3)
  (2,4)
  (2,5)
  (2,6)
  (2,7)
  (2,8)
  (2,9)
  (2,10)
  (2,11)
  (2,12)
  (2,13)
  (2,14)
  (2,15)
  (2,16)
  (2,17)
  (2,18)
  (2,19)
  (2,20)
  (2,21)
  (2,22)
  (3,1)
  (3,2)
  (3,3)
  (3,4)
  (3,5)
  (3,6)
  (3,7)
  (3,8)
  (3,9)
  (3,10)
  (3,11)
  (3,12)
  (3,13)
  (3,14)
  (3,15)
  (3,16)
  (3,17)
  (3,18)
  (3,19)
  (3,20)
  (3,21)
  (3,22)
  (3,23)
  (3,24)
  (3,25)
  (3,26)
  (3,27)
  (3,28)
  (3,29)
  (3,30)
  (3,31)
  (3,32)
  (4,1)
  (4,2)
  (4,3)
  (4,4)
  (4,5)
  (4,6)
  (4,7)
  (4,8)
  (4,9)
  (4,10)
  (4,11)
  (4,12)
  (4,13)
  (4,14)
  (4,15)
  (4,16)
  (4,17)
  (4,18)
  (4,19)
  (4,20)
  (4,21)
  (4,22)
  (4,23)
  (4,24)
  (4,25)
  (5,1)
  (5,2)
  (5,3)
  (5,4)
  (5,5)
  (5,6)
  (5,7)
  (5,8)
  (5,9)
  (5,10)
  (5,11)
  (5,12)
  (5,13)
  (5,14)
  (6,1)
  (6,2)
  (6,3)
  (6,4)
  (6,5)
  (6,6)
  (6,7)
};
\end{axis}
\end{tikzpicture}
\end{center}
\begin{firsttakebox}{FIRST TAKE --- before the video (5 min)}
Which claim seems better supported right now? Cite one detail from the sequence or dotplot.
\workspace{0.9in}
\end{firsttakebox}

\begin{callout}{\IconBook\ MODEL THE PATTERN, NOT THE STORY (keep on the page)}{calloutgreen}
Random processes can produce visible patterns, including streaks. To judge a pattern under a proposed chance model,
make each simulation trial match the original process, choose a statistic that measures how extreme the pattern is,
and count results \textbf{at least as extreme} as the observation. Repeat many independent trials; the successful-trial
proportion estimates how often chance alone would produce such a result. An unusual result is evidence against a model,
not proof of what caused the departure.
\end{callout}

\newpage
\textbf{Finish after the video:}\par\smallskip

\dokbadge{1}~\textbf{(a)} State the longest run in the observed sequence and the most common longest-run value among the simulated sets.\par
\workspace{0.7in}

\dokbadge{2}~\textbf{(b)} Use the dotplot to estimate the probability that 10 fair-coin flips have a longest run at least as long as the observed run. Interpret the estimate in context.\par
\workspace{1.3in}

\dokbadge{3}~\textbf{(c)} Which claim is better supported by this evidence? Cite the simulation result that settles your choice. Then describe a larger simulation that would more precisely settle the disagreement: state what one trial is, what to record and count, how many trials to use, and what result would make you doubt that the coin is fair.\par
\workspace{2.0in}

\begin{sentenceframebox}\raggedright\textbf{Frame:} Claim \blank[0.35in] is better supported because \blank[2.7in].\end{sentenceframebox}

\begin{sentenceframebox}\raggedright\textbf{Frame:} One trial is \blank[1.8in]; I would record \blank[1.8in] and doubt fairness if \blank[1.8in].\end{sentenceframebox}

\begin{summaryexitbox}{EXIT --- turn this in}
One thing the video changed about my first take: \hrulefill\par
\workspace{0.4in}
\end{summaryexitbox}

\end{document}
